\documentclass{article}
\usepackage[utf8]{inputenc}
\usepackage[T1]{fontenc}
\usepackage[margin=2cm]{geometry}
\usepackage{hyperref}
\sloppy
\begin{document}
\section*{Historical source}
Encyclopaedia Britannica scan, source row 1680. Parse 3.5 extracted the text. Super 3.5 VL corrected it from an enlarged crop of each block. It can still contain recognition errors.
Page headers, footers, and page numbers are removed.

\subsection*{CARPENTRY (E0147)}
\textit{Continues an entry from an earlier page.}

Carpentry.

try. Aσ, and 'Ag. But because As and cE are equal and parallel, and Er and fi are also parallel, as also er and fp, it is evident, that if is equal to rE, or to oF, and iA is equal to re, or to Gg. Therefore the four forces Ag, Aσ, Ap, Ai, are equal to AG and AS. Therefore AG is the compression of the beam AD, or the force pressing it on D, and AS is the force pressing it on the beam HD. The proportion of these pressures, therefore, is not affected by the form of the joint.

This remark is important; for many carpenters think the form and direction of the butting joint of great importance; and even the theorist, by not prosecuting the general principle through all its consequences, may be led into an error. The form of the joint is of no importance, in as far as it affects the strains in the direction of the beams; but it is often of great consequence, in respect to its own firmness, and the effect it may have in bruising the piece on which it acts, or being crippled by it.

Origin of the Strain on a Tbe beam.

The same compression of AB, and the same thrust on the point D by the intervention of AD, will obtain, in whatever way the original pressure on the end A is produced. Thus, supposing that a cord is made fast at A, and pulled in the direction AE, and with the same force, the beam AD will be equally compressed, and the prop D must react with the same force.

But it often happens that the obliquity of the pressure on AD, instead of compressing it, stretches it; and we desire to know what tension it sustains. Of this we have a familiar example in a common roof. Let the two rafters AC, AD (fig. 12.), press on the tie-beam DC. We may suppose the whole weight to press vertically on the ridge A, as if a weight B were hung on there. (See Note CC.) We may represent this weight by the portion Ab of the vertical or plumb line, intercepted between the ridge and the beam. Then drawing bf and bg parallel to AD and AC, Ag and As will represent the pressures on AC and AD. Produce AC till CH be equal to As. The point C is forced out in this direction, and with a force represented by this line. As this force is not perpendicularly across the beam, it evidently stretches it; and this extending force must be withstood by an equal force pulling it in the opposite direction. This must arise from a similar oblique thrust of the opposite rafter on the other end D. We concern ourselves only with this extension at present; but we see that the cohesion of the beam does nothing but supply the balance to the extending forces. It must still be supported externally, that it may resist, and by resisting obliquely, be stretched. The points C and D are supported on the walls, which they press in the directions CK and DO, parallel to Ab. If we draw HK parallel to DC, and HI parallel to CK (that is to Ab), meeting DC produced in I, it follows from the composition of forces, that the point C would be supported by the two forces KC and IC. In like manner, making DN = Ag, and completing the parallelogram DMNO, the point D would be supported by the forces OD and MD. If we draw go and fk parallel to DC, it is plain that they are equal to NO and CI, while Ao and Ak are equal to

DO and CK, and Ab is equal to the sum of DO and Carpentry. CK (because it is equal to Ao + Ak). The weight of the roof is equal to its vertical pressure on the walls.

Thus we see, that while a pressure on A, in the direction AB, produces the strains AS and AG, on the pieces AC and AD, it also excites a strain CI or DM in the piece DC. And this completes the mechanism of a frame; for all derive their efficacy from the triangles of which they are composed, as will appear more clearly as we proceed.

But there is more to be learned from this. The consideration of the strains on the two pieces AD and AC, by the action of a force at A, only showed them as the means of propagating the same strains in their own direction to the points of support. But, by adding the strains exerted in DC, we see that the frame becomes an intermediate, by which exertions may be made on other bodies, in certain directions and proportions; so that this frame may become part of a more complicated one, and, as it were, an element of its constitution. It is worth while to ascertain the proportion of the pressures CK and DO, which are thus exerted on the walls. The similarity of triangles gives the following analogies:

DO : DM : OK : DI : CM : = OO A Where Do : CK - CE: All

Or, the pressures on the points C and D, in the direction of the straining force Ab, are reciprocally proportional to the portions of DC intercepted by Ab.

Also, since Ab is = DO + CK, we have
Ab : CK = Cb + bD (or CD) : bD, and
Ab : DO = CD : bC.

In general, any two of the three parallel forces AB, DO, CK, are to each other in the reciprocal proportion of the parts of CD, intercepted between their directions and the direction of the third.

And this explains a still more important office of the frame A D C. If one of the points, such as D, be supported, an external power acting at A, in the direction A b, and with an intensity which may be measured by A b, may be set in equilibrio, with another acting at C, in the direction C L, opposite to C K or A b, and with an intensity represented by C K: for since the pressure C H is partly withstood by the force I C, or the firmness of the beam D C supported at D, the force K C will complete the balance. When we do not attend to the support at D, we conceive the force A b to be balanced by K C, or K C to be balanced by A b. And, in like manner, we may neglect the support or force acting at A, and consider the force D O as balanced by C K.

Thus our frame becomes a lever, and we are able to trace the interior mechanical procedure which a Lever gives it its efficacy: it is by the intervention of the forces of cohesion, which connect the points to which the external forces are applied with the supported point or fulcrum, and with each other.

ported point of fulcrum, and with each other.
These strains or pressures Ab, DO, and CK, not
being in the directions of the beams, may be called
transverse. We see that by their means a frame of


\section*{Source record}
\begin{description}
\item[Scan ID] \texttt{britannica:309f765b1ecc83e0a7f7762e44b4a927a7555443:train:1680}
\item[Image file] \texttt{page-1680-fbb4a5fde0c2.jpg}
\item[Image SHA-256] \texttt{fbb4a5fde0c2b233b6df827e5998ef2fc15cd2bc73b4eb99b93d91db83199d1b}
\item[Dataset] NationalLibraryOfScotland/encyclopaedia\_britannica\_illustrated at revision \texttt{309f765b1ecc83e0a7f7762e44b4a927a7555443} (CC0-1.0)
\item[Extraction] model \texttt{nvidia/nemotron-3.5-parse}, raw reply \texttt{extract/raw/page-1680-fbb4a5fde0c2.json}
\item[Raw reply SHA-256] \texttt{874ef4a847c3c35b6d1c6fec79439e161eeefcfbb1daafe8853a24cd7f3aad09}
\item[Correction] model \texttt{nvidia/Nemotron-3.5-Super-VL}, block crops, raw replies \texttt{correct/raw/crop\_nothink/page-1680-fbb4a5fde0c2-b*.json}; parse: 1, parse\_fallback: 1, super\_vl: 13
\item[Running header] C ARPENTRY. 631
\end{description}
\end{document}
